\documentclass[12pt,a4paper]{article}
\usepackage[utf8]{inputenc}
\usepackage[T1]{fontenc}
\usepackage{lmodern}
\usepackage[margin=24mm]{geometry}
\usepackage{xcolor,parskip}
\usepackage[colorlinks=true,urlcolor=teal]{hyperref}
\definecolor{ink}{HTML}{152D40}
\color{ink}
\setlength{\parindent}{0pt}
\pagestyle{plain}
\begin{document}

{\small Speaking script · 10-minute talk · 1--2 October 2026}\par
\section*{1. Rotating black holes with AI agents}
\textbf{0:00--0:40 (40 seconds)}\par
{\itshape Introduce the question. Do not read the affiliations aloud.}\par\medskip
We used AI agents to work on a concrete numerical relativity problem: rotating black holes in five-dimensional Einstein--Gauss--Bonnet gravity. The project had two stages. First, reproduce a known numerical family. Then use that solver to ask how a higher-curvature term changes the extremal mass at fixed angular momentum. I will show the evidence, the remaining limitations, and what the agents actually cost. Claude and ChatGPT, used through Codex, assisted with the work. The four human authors are responsible for the scientific claims.
\newpage
{\small Speaking script · 10-minute talk · 1--2 October 2026}\par
\section*{2. The first prompts}
\textbf{0:40--1:50 (70 seconds)}\par
{\itshape Point to the original sentence, then explain how the scope narrowed.}\par\medskip
The project began on September fourth with a deliberately broad request: use agents to implement numerical codes for rotating black holes above four dimensions. But the first prompt also set an order: reproduce known solutions, verify the codes, and only then look for something new. The sentence on the slide is a literal excerpt from that prompt.

The next prompt supplied a literature list, including the 2010 paper by Brihaye and collaborators. A third asked the agent to follow the course requirements. So the starting point was a research objective and reference material, rather than a fully specified algorithm.

The working problem became asymptotically flat black holes in five dimensions with two equal angular momenta. That symmetry reduces the field equations to radial ordinary differential equations. It made the calculation manageable while retaining a genuinely nonlinear rotating problem.
\newpage
{\small Speaking script · 10-minute talk · 1--2 October 2026}\par
\section*{3. Reproducing the profiles in 1010.0860}
\textbf{1:50--3:20 (90 seconds)}\par
{\itshape Point first to the profiles, then spend time on the difference panel.}\par\medskip
This is the reproduction test against Figure one of the 2010 paper. The four curves are metric functions. Lines show the project calculation, and circles show vertices extracted from the published figure's PostScript source. The axes and conventions match, including a factor of four between the two definitions of the Gauss--Bonnet coupling.

On the left, the agreement looks very good. The difference panel on the right is essential: the largest absolute mismatch is about zero point zero one four, well above the figure's plotting resolution. We should therefore describe this as a reproduction with a quantified discrepancy, rather than perfect agreement.

A near-Einstein control also finds a discrepancy between the published plotted curve and the analytic solution. That limits what this figure can establish, but it does not identify the cause. A separate outer-ergosurface comparison also remains unresolved.

The stronger checks use analytic limits, the full field equations, and thermodynamic identities. This was an early lesson for the agents: a convincing overlay is only the beginning of validation.
\newpage
{\small Speaking script · 10-minute talk · 1--2 October 2026}\par
\section*{4. Measuring the shift of extremality}
\textbf{3:20--4:40 (80 seconds)}\par
{\itshape Explain Psi in words. The response equation is supporting detail.}\par\medskip
The new question concerns the mass of the extremal branch at fixed individual angular momentum. We add the Gauss--Bonnet invariant to the Einstein action, with coupling alpha. In the extended first law, Psi measures the mass response to alpha at fixed entropy and angular momenta. The factor of two counts the two equal spins.

On a regular extremal branch, the temperature vanishes and Psi becomes the derivative of the extremal mass with respect to alpha. That is the quantity we want.

The project solves for four radial functions using spectral collocation and Newton iteration. It then differentiates the discrete equations themselves. The equation on the right gives the response through one linear solve, without choosing a finite-difference step in alpha. This is exact differentiation of the discrete problem, not an exact continuum solution.

The archive contains 356 accepted states at thirteen couplings. The zero-temperature values come from extrapolating those finite-temperature solutions.
\newpage
{\small Speaking script · 10-minute talk · 1--2 October 2026}\par
\section*{5. The slope stays positive and turns upward}
\textbf{4:40--6:20 (100 seconds)}\par
{\itshape Trace the black points from left to right. Pause at the valley.}\par\medskip
Here is the main result. The horizontal axis is the coupling in units set by angular momentum. The vertical axis is the derivative of the extremal mass at fixed angular momentum. The green star is the known perturbative result: the initial slope is pi. That value comes from the literature, and the project recovers it.

As the coupling increases, the black points fall to about one point zero seven, near y equal to zero point one seven five. They then rise to about one point four four at the largest sampled coupling. Every central value remains positive. So the extremal mass keeps increasing, but its rate of increase is not constant.

The bars measure how much the extrapolated answer changes when we change the fitting prescription. They are not confidence intervals and do not include every possible systematic error. The broad fall and rise are much larger than these fit sensitivities, but the precise location of the minimum between samples remains unresolved.

The open squares are slopes between neighboring mass points. They support the broad trend, although they share the same underlying numerical solutions. They are a consistency check, not a fully independent measurement.
\newpage
{\small Speaking script · 10-minute talk · 1--2 October 2026}\par
\section*{6. The extremal mass grows more slowly than linear order}
\textbf{6:20--7:35 (75 seconds)}\par
{\itshape Contrast the black mass curve with the dashed first-order prediction.}\par\medskip
This is the same result expressed as a mass curve. At fixed angular momentum, the extremal mass grows by about thirty-one percent over the sampled range. The dashed line uses only the first-order correction. It increasingly overestimates the mass, because the measured slope falls below pi.

The contribution here is the finite-coupling derivative and its numerical characterization. The rotating family and its extremal branch were already known. We should not present them as new solutions.

There are also clear boundaries to the claim. Every extremal point shown is an extrapolation. The project does not establish that this branch is unique, stable, or globally lowest in mass. And the calculation treats classical Einstein--Gauss--Bonnet gravity as the theory. In a general effective field theory, additional higher-curvature terms can matter at the same order. So this curve is not a universal prediction for all corrections to Einstein gravity.
\newpage
{\small Speaking script · 10-minute talk · 1--2 October 2026}\par
\section*{7. Recorded time and token use}
\textbf{7:35--9:05 (90 seconds)}\par
{\itshape Emphasize what the numbers count before comparing the work blocks.}\par\medskip
The recovered usage report covers September fourth through September twenty-eighth. It records about thirty-seven hours of active agent time, one point four two billion processed tokens, and five point nine four million output tokens.

The large processed-token number mostly counts repeated reads of cached context. It is not that much newly written text, and it is not a monetary cost. The two tools also account for tokens differently, so their totals should not be used as an efficiency ranking.

The most expensive block was building the rotating Gauss--Bonnet equations and solver: eleven point two hours. The response calculation, manuscript and review took eight point seven hours. The first attempt at the conjugate potential took four point five hours. The solver block consumed roughly five times the processed tokens of the earlier analytic and static benchmarks together.

These are active-session estimates using a twenty-minute inactivity cutoff. They exclude human work and unattended numerical runs. We therefore do not have a defensible total CPU or GPU computation time. The cost record is useful precisely because we keep that distinction.
\newpage
{\small Speaking script · 10-minute talk · 1--2 October 2026}\par
\section*{8. What we learned about agents in research}
\textbf{9:05--10:00 (55 seconds)}\par
{\itshape End on the physics result and the next decisive check.}\par\medskip
The agents helped turn a research question into a solver, a dataset, and a manuscript. The useful discipline was to connect claims to checks: analytic benchmarks, independent evaluations of the equations, and saved records of failures. The expensive parts included unsuccessful methods and corrections, not just writing the final code.

Human judgment still determined the scope and what the evidence justified. Agreement between several diagnostics is less reassuring when they share the same source of error.

The result to take away is a positive, non-monotonic extremal slope on the sampled equal-spin branch. The next decisive check is an independent construction of extremal solutions and a comparison of the slope. That would address the main remaining dependence on finite-temperature extrapolation.
\newpage\section*{Sources and preparation notes}
Not part of the spoken ten minutes. Timing assumes a deliberate pace with pauses to point at the figures. Rehearse once to match your own delivery.
\subsection*{Scientific sources}
Slides 3--6 summarize short-summary/summary.tex and manuscript/main.tex, using their existing figures and generated numbers. No new black-hole solutions were computed for this presentation. The response identity is rechecked algebraically in validate.py. The first law and analytic benchmarks are established inputs. The finite-coupling numerical response is the project result.
\subsection*{Papers already stored in the project}
Brihaye, Kleihaus, Kunz and Radu, arXiv:1010.0860, supplies the numerical family and Figure 1 comparison. Ma, Li and Lü, arXiv:2009.00015, supplies the first-order extremal mass correction. Kleihaus, Kunz and Radu, arXiv:2303.12471, already constructs the extremal branch directly. All three source PDFs are in papers/.
\subsection*{Figures}
The comparison and slope figures are unchanged project figures. Slide 6 shows only the left panel of fig-massext.pdf, with all axes of that panel preserved. The PDF slides retain vector figures. The HTML embeds a vector rendering of every slide and a searchable speaking script. No external fonts, scripts or images are required.
\subsection*{Prompt provenance}
Slide 2 quotes entry 001, dated 4 September 2026, in prompts/prompt-log-cronologico.md. Entries 002 and 003 supply the literature list and course-compliance request. The English emphasis line is a paraphrase, not another historical prompt.
\subsection*{Usage provenance and limitations}
Slide 7 transcribes prompts/uso-de-tiempo-y-tokens.md. Exact totals: 1,415,508,752 processed tokens and 5,942,462 output tokens. Claude Code cache reads alone account for 1,301,219,499 processed tokens. The report totals active time as 36.8 h, while its milestone table totals 36.6 h; this unresolved 0.2 h discrepancy is retained here rather than silently reconciled. Varying the inactivity cutoff from 10 to 30 minutes changes active time from 30.2 to 39.3 h. The report excludes human work, unattended numerical runs and missing sessions. These historical totals exclude preparation of this presentation. They were not recomputed from the original private session logs.
\newpage
\subsection*{AI acknowledgment}
Suggested wording: “AI assistance: Claude (Anthropic) and ChatGPT / Codex (OpenAI), for code development, calculations, checks and drafting. The human authors take responsibility for the scientific claims.” The stored workflow uses Claude Code and Codex. The acknowledgment distinguishes assistance from authorship without implying that the tools verified their own output independently.
\subsection*{Affiliations and event}
Web check: 29 September 2026. Cecilia Bejarano and Andrés Goya: IAFE, CONICET--UBA. David Blanco and Guillem Pérez-Nadal: Departamento de Física, FCEN, UBA and IFIBA, CONICET--UBA. Institutional listings support the appointments; department and institute names are expanded on the title slide. Course title, leaders, venue and 1--2 October dates follow the presentation request. The year is 2026, consistent with the project history and current date.
\subsection*{Affiliation links}
\href{https://www.iafe.uba.ar/relatividad/gravity.html}{IAFE Gravitation and Cosmology group}. Lists Cecilia Bejarano and Andrés Goya as permanent researchers, and David Blanco as an external collaborator at IFIBA, CONICET--UBA.\par
\href{https://wp.df.uba.ar/hepth/people/}{UBA high-energy theory group}. Lists David Blanco, Andrés Goya and Guillem Pérez-Nadal among faculty.\par
\href{https://df.uba.ar/es/investigadores1/investigadores-y-becarios/miembro/75-Guillem_Perez_Nadal}{UBA researcher directory: Guillem Pérez-Nadal}. Lists a UBA faculty appointment and CONICET researcher status.\par
\href{https://www.df.uba.ar/attachments/article/8731/Newsletter%20mayo%202025.pdf}{UBA department newsletter, May 2025}. Identifies Blanco, Goya and Pérez-Nadal among UBA--CONICET researchers organizing the school.\par
\end{document}